\section{2교시 (90분) — 서비스 전환: ITIL · SAC · SLA 세 층 · error budget}
% =====================================================================
% 2교시 — 교재 제3장, 워크북 §2(⑰ 서비스 전환 패키지). 모든 숫자는 Day4 워크북 완성 예시본(정본).
% =====================================================================
\begin{frame}{2교시 — 이 교시에서 배우는 것}
\begin{itemize}
\item 약속(SLA)은 그것을 받치는 약속(OLA·UC)의 합보다 강할 수 없다
\item SAC는 거부권이 실재할 때만 장치가 된다 — 판정자 박준영
\item ELS는 기간·종료 기준·인력 축소가 연동된 8주
\item 99.5\%는 219.2분이고, 그 예산이 릴리스 횟수를 정한다
\item 인시던트는 복구로, 문제는 원인 제거로 닫는다 — KEDB 6건
\end{itemize}
\keymsg{2교시의 한 문장 — 약속과 그 약속을 받치는 것을 분리하지 않는 것이 이 장의 모든 오류다}
\note{교재 제3장 도입과 3.11절. 1교시가 "되돌릴 수 없는 시각"(컷오버)이었다면 2교시는 그 다음 날부터 운영 조직이 단독으로 서비스를 지원할 수 있을 때까지의 구간이다. 워크북 §2 ⑰ 서비스 전환 패키지(SAC·ELS·SLA/OLA/UC 3층·error budget policy·KEDB)가 이 교시의 산출물이다. 순서는 어휘(ITIL 판본 8분) → 인수(SAC) → 기간(ELS) → 회의체(CAB) → 약속의 산술(SLA 3층·error budget) → 사건(KEDB·개인정보) → 사람(채택) → 절차의 부작용(D$-$2)이다. 예상 소요 3분.}
\end{frame}

\begin{frame}{ITIL 판본을 정직하게 — v5가 현행, 이 교재는 practice 어휘로 고정}
\fig[0.58]{../그림/PM_Day4/fig_3_1_itil_vocab.pdf}
\figcap{그림 3-1. ITIL Version 5(2026-02-12) 상위 구조 — ITIL Value System · Product and Service Lifecycle 8활동. 이 장은 상위 구조 명칭을 쓰지 않고 practice 명칭 9개만 쓴다}
\keymsg{"ITIL 4가 최신"은 2026-02-12 이후 틀린 문장 — 바뀐 것은 상위 구조·6\ra{}8활동, practice 이름은 그대로}
\note{교재 3.1절. 판본 논쟁을 8분 안에 끝내기 위한 프레임이다. 사실 — 2026-02-12 PeopleCert가 ITIL v5 Foundation 출시, 상위 모듈 3\til{}4월 순차. 바뀐 것 둘: Service Value System \ra{} ITIL Value System, Service Value Chain 6활동 \ra{} Product and Service Lifecycle 8활동. 유지된 것: 34 practices·7 원칙·4 차원. 이 교재가 ITIL 4 어휘(Change Enablement·Service Validation and Testing·Release·Deployment·Incident·Problem·Service Level Management·ELS·KEDB)로 고정하는 이유 셋 — 회사 운영 규정과 벤더 계약이 그 어휘 / practice 이름은 판본 무관 / 다투는 것은 어휘가 아니라 판단. 회사에 돌아가 규정을 개정할 때 바꿀 것은 상위 구조 명칭과 6\ra{}8 매핑뿐이다. 서비스 전환(Transition)은 8활동 중 Transition과 Operate의 경계이고 ELS는 그 경계 위의 기간이다. 예상 소요 6분.}
\end{frame}

\begin{frame}{SAC — 인수 거부권이 실재해야 문서가 장치가 된다}
\begin{itemize}
\item 판정자 = 인수자(운영 조직, 박준영) — 프로젝트가 판정하면 자기 평가
\item 거부권은 헌장 §14 ①에 걸려 있어야 실재한다
\item 조건부 충족 경로 — 없으면 1항목 때문에 전체 거부 또는 거짓 충족
\item 시점 — G1 초안 \ra{} 월 검토 \ra{} G3 사전 점검 \ra{} G4 판정
\item 항목 12 · 각 항목에 판정자와 증빙 열
\end{itemize}
\keymsg{프리모템 12번("미해소 23건으로 인수 거부")과의 차이는 항목 수가 아니라 G3 사전 점검의 유무다}
\note{교재 3.2절 개념, 워크북 §2 항목 1\til{}2. SAC(Service Acceptance Criteria)는 ITIL Service Validation and Testing이 원전, PRINCE2 7 Quality practice의 acceptance criteria \ra{} acceptance record 사슬. 항목 10\til{}12개 — 운영 문서·모니터링·백업/복구·보안·교육·Known Error·SLA/OLA/UC 서명·인시던트/문제/변경 절차·소스/형상·인터페이스 운영·라이선스/계약 이관·잔여 결함/ELS 종료. 장치가 되는 조건은 하나 — 판정자가 인수자이고 그에게 거부권이 실재. 거부권이 실재하면 두 번째 조건이 따라온다 — 조건부 충족의 경로("미해소 0 또는 기한·책임자 서명"). 조건부가 없으면 인수자는 1항목 때문에 전체를 거부하거나(하자보수기 지연·유보금·팀 해산 지연) 거짓으로 충족을 서명한다. 예상 소요 6분.}
\end{frame}

\begin{frame}{SAC 12항목 — G3 사전 점검(미해소 7)에서 G4 판정(충족 11·조건부 1)까지}
\fig[0.58]{../그림/PM_Day4/fig_3_2_sac_progress.pdf}
\figcap{그림 3-2. 01-28 미해소 7 \ra{} 14주 해소 이벤트(DR 시험 04-10 · 임계값 12\ra{}24 · 데스크 교육 11/11 · CAB 8회 · UC 재협상 · 문서 12종 05-14) \ra{} 05-21 판정. 조건부 1(항목 11 아카이브 접근 절차, 재경팀 06-15) = 이관 목록 1번}
\keymsg{사전 점검이 있었기 때문에 미해소 7은 작업 목록이 되었고, 조건부 1은 이관 목록 1번이 되었다}
\note{교재 3.2절 온담 사례, 워크북 §2 항목 2·§1.5 항목 8. SAC v1.0은 G1(2026-04-24) 초안 합의(RACI 22행), 월 검토 9회, 사전 점검 2027-01-28, 판정 05-21, 판정권자 박준영. 01-28 미해소 7 — 운영 문서 8/12, 모니터링 임계값 12/24, DR 시험 미실시, 데스크 교육 6/11, SLA UC 3PL 없음·수행사 UC 초안, CAB 미실시, 라이선스·계약 이관 미해소. 조건부 1은 항목 11의 아카이브 접근 절차 — 문서와 운영팀 검토는 끝났으나 재경팀의 세무 보존 요건 정합 확인이 06-15에야 나온다. 이것을 "미해소"로 두면 G4가 밀리고 "충족"으로 두면 거짓이다 — 그래서 조건부 경로가 필요하고, 조건부는 기한(06-15)·책임자 서명을 갖고 이관 목록 1번이 된다. 항목 7이 "SLA 서명"이 아니라 "3층 전부 서명"인 이유는 뒤의 SLA 프레임에서. 예상 소요 6분.}
\end{frame}

\begin{frame}{ELS 8주 — 종료 조건이 없는 하이퍼케어는 상주가 된다}
\fig[0.58]{../그림/PM_Day4/fig_3_3_els_curves.pdf}
\figcap{그림 3-3. 2027-03-01\til{}04-23 · 오류율 2.4\ra{}0.4\% · 콜 1,840\ra{}450 · P1 2(1주차)·P2 9 · 상주 20\ra{}15\ra{}10\ra{}6\ra{}4 · 5주차 재상승은 2차 컷오버 유예 조건 안 · 7·8주 2주 연속 충족 \ra{} 04-23 종료}
\keymsg{ELS의 세 요소 — 기간 · "n주 연속"의 종료 기준 · 지표와 연동된 인력 축소. 하나라도 빠지면 상주는 하자보수가 된다}
\note{교재 3.3절, 워크북 §2 항목 5. 종료 기준 다섯 — ① 2주 연속 P1 0, ② P2 주 2건 이하, ③ 매장 주문 오류율 0.5\% 이하(O-13 1.8\% \ra{} 목표), ④ 데스크 콜 480/주 이하, ⑤ 심각도 2 backlog 0. 인시던트 140 = P1 2·P2 9·P3 41·P4 88, 핫픽스 11, 판정 박준영. 상주 축소가 달력이 아니라 지표에 연동되어 있다 — 3주차 15는 2주차 판정(심각도 2 이상 4 $\le$ 5) 후, 6주차 10은 오류율 0.5\% 도달 후, 8주차 6은 7주차 기준 충족 후. 5주차 재상승(P2 2·오류율 0.8\%, P2-07 폴백 배치 84분)은 "2차 후 1주는 기준 적용 유예"가 미리 적혀 있었기 때문에 ELS 연장 토론이 없었다. 가용성 산정 — 3월 다운 57분은 P1-01만 산입, P1-02(포인트 이중차감)는 데이터 손상이지 다운이 아니다. 산입 규칙이 미리 없으면 "왜 안 넣었나"가 된다. 예상 소요 7분.}
\end{frame}

\begin{frame}{Change Enablement · Release · Deployment — 셋을 구분해야 CAB가 정치가 되지 않는다}
\begin{tbl}[\scriptsize]{L{0.2\textwidth}L{0.24\textwidth}L{0.48\textwidth}}
practice & 답하는 질문 & 요소 \\ \midrule
Change Enablement & 이 변경을 해도 되는가(승인) & standard·normal·emergency · change authority · CAB/ECAB · 변경 일정 \\
Release Management & 무엇을 언제 묶어 내보내는가(구성) & 릴리스 단위 · 릴리스 일정 · 릴리스 노트 \\
Deployment Management & 어떻게 환경에 옮기는가(실행) & big-bang·phased·pull·continuous · 배포 절차 · 롤백(런북) \\
\end{tbl}
\begin{itemize}
\item 섞이면 CAB가 "이 기능을 넣을 것인가"와 "이번 주 배포해도 되는가"를 한 회의에서 다룬다
\item 온담 — CCB(정미란) 05-28 해산 · CAB(박준영, 수요일) 03-17 첫 회, 8회 = SAC 항목 8 증빙
\end{itemize}
\keymsg{무엇을 넣을지는 릴리스 계획(현업), 언제 어떻게 내보낼지는 CAB(운영 위험 · error budget), 표준 변경은 CAB에 오지 않는다}
\note{교재 3.4절, 워크북 §2 항목 6. 세 practice는 다른 질문에 답한다. 섞이면 제품 결정(현업 우세)과 운영 위험 결정(운영 우세)이 한 회의에 올라 정치가 되고, 나누면 계산이 된다 — 그 계산의 환율이 error budget이다. 운영 이관은 변경 권한의 이전이기도 하다 — 프로젝트 중 CCB(정미란 위원장), 이관 후 CAB(박준영 위원장). 겹치는 기간(ELS)의 규칙: CAB 첫 회 03-17\til{}CCB 해산 05-28 사이 프로젝트 변경(CR-14 잔여·R5 릴리스)은 CCB, 운영 변경(핫픽스·파라미터)은 CAB. 이관 3번(22:00 운용 전환)은 CCB 01-21에서 "P2 안정화 후 CAB 재상정"으로 넘어간 것 — 프로젝트 변경이 운영 변경으로 권한 이전된 실물. 표준 변경 목록이 없으면 파라미터 하나도 2주 걸린다. 예상 소요 6분.}
\end{frame}

\begin{frame}{SLA/OLA/UC 3층 — 한 층이 비면 SLA는 협상이 아니라 산술로 불가능하다}
\fig[0.56]{../그림/PM_Day4/fig_3_4_sla_chain.pdf}
\figcap{그림 3-4. 항목 ② P1 복구 2h — 사슬 15+30+30+240 = 315분 $>$ 120분(UC 초안 불가) \ra{} UC 재협상(착수 30·복구 90, 연 +0.6억) 165분 \ra{} 병렬 착수 $\le$ 120 성립. 항목 ⑥ 출고 확정 D+1 06:00 — 3PL UC 없음 \ra{} 최선 노력 강등}
\keymsg{SLA 복구 시간 $\ge$ 받치는 OLA·UC 사슬의 합 — 강등은 패배가 아니라 정직한 경계 표시이자 재협상 안건 목록이다}
\note{교재 3.5절, 워크북 §2 항목 7(3층 매핑표). SLA = IT본부 \ra{} 현업, OLA = IT본부 내부 팀 간(서비스운영 \ra{} 인프라·보안·개발운영), UC = 외부 공급자(수행사 유지보수·CSP·앱 벤더·3PL·API GW 벤더). 온담 SLA 7·OLA 3·UC 5, 서명 05-14. ②는 G3 사전 점검(01-28)에서 315분이 발견되어 05-14 전에 재협상이 끝났다 — G4에서 처음 발견되었다면 SLA는 발효되지 못한 채 이관되었을 것. ⑥은 출고 확정이 대성로지스 WMS의 작업이고 계약은 물류본부 소관(P1 예산 밖), 준실시간 연동은 거부(R-03·IS-05) — UC 자체가 없으므로 어떤 OLA로도 성립하지 않는다. 2027-06 3PL 재협상 후 재상정 = 이관 4번(한동석, 06-30). 나머지: ① 가용성 99.5\%(CSP 99.9\%), ③ P2 8h(330 $\le$ 480), ⑦ 앱 주문 연동 4h(210 $\le$ 240, R-18 월 3,600만 원, 이관 5번). 예상 소요 8분.}
\end{frame}

\begin{frame}{Error budget — 99.5\%를 3시간 39분으로, 다시 릴리스 속도로}
\begin{itemize}
\item 평균 달 = 365.25/12 $\times$ 1,440 = 43,830분 · 0.5\% = \textbf{219.2분 = 3h 39m}
\item 계획 정지(CAB 승인 창·2차 컷오버 8h)는 제외 — 산입하면 3월 98.7\%
\item 비계획 유보 60 : 릴리스 몫 40 = 131.5분 : \textbf{87.7분}
\item 릴리스당 기대 다운 = 변경 실패율 $\times$ 복구 시간(DORA)
\item policy — 유보 초과 시 표준 릴리스 동결 · 릴리스 몫 소진 시 CAB 승인만 · 결정자 박준영
\end{itemize}
\keymsg{같은 99.5\%를 운영은 "지킬 숫자"로, 프로젝트는 "넘지 않을 숫자"로 읽는다 — 예산으로 바꾸면 토론이 계산이 된다}
\note{교재 3.6절 개념, 워크북 §2 항목 3(error budget policy). SRE(Beyer 외, 2016, 3·4장). 28일 달 202분·30일 216분·31일 223분. 예산은 장애(비계획)에만 쓰는 것이 아니라 릴리스에도 쓴다 — 릴리스 한 번은 기대값으로 실패율 $\times$ 복구 시간만큼 먹는다. policy가 없으면 예산이 소진되어도 아무 일도 안 일어나고, 그러면 예산이 아니다. 계획 정지 제외 규칙이 미리 없으면 2차 컷오버 8시간이 들어가 헌장 §14 ④가 흔들린다. ELS 실적: 3월 다운 57분(P1-01) = 223분의 26\%, 유보의 43\% / 4월 0분. R5 릴리스(05-07)와 핫픽스 11건이 이 예산 안에서 결정되었다. 다음 프레임에서 환전 계산. 예상 소요 5분.}
\end{frame}

\begin{frame}{환전 계산 — 릴리스 속도를 여섯 배로 만드는 것은 기능이 아니라 복구 시간이다}
\fig[0.5]{../그림/PM_Day4/fig_3_5_error_budget.pdf}
\figcap{그림 3-5. 219.2분 도넛(유보 60\% / 릴리스 40\% = 87.7, 3월 실측 57분) · 릴리스당 기대 다운 대 월 릴리스 가능 횟수}
\begin{tbl}[\scriptsize]{L{0.34\textwidth}L{0.14\textwidth}L{0.2\textwidth}L{0.2\textwidth}}
릴리스당 기대 다운 & 분 & 월 가능(87.7 $\div$) & 주 환산 \\ \midrule
DORA 실측 12\% $\times$ 240분 & 28.8 & 3.0회 & 0.7회 — 주 1회 불가 \\
중간(하자보수기) 10\% $\times$ 90분 & 9.0 & 9.7회 & 2.2회 — 표준 주 1회 + 핫픽스 \\
ELS 목표 10\% $\times$ 45분 & 4.5 & 19.5회 & 4.5회 \\
\end{tbl}
\keymsg{240분 \ra{} 45분은 KEDB(우회 준비된 오류) · 회귀 자동화(R-19 0.25억) · 롤백 스크립트의 문제 — 안정과 속도의 환율}
\note{교재 3.6절 온담 사례, 워크북 §2 항목 3. DORA 실측(2026-09, Day3 2.7절)으로는 월 3회가 한계이고 ELS 목표로는 19회. 차이를 만드는 변수는 실패율(12\ra{}10\%)보다 복구 시간(240\ra{}45분)이다. 복구 45분은 기능 팀이 아니라 KEDB·회귀 자동화·롤백 스크립트의 문제 — 그래서 KEDB 6건이 SAC 항목 6이고 R-19 회귀 도구 0.25억이 릴리스 속도 투자다. 이 환율이 있으면 CAB는 "위험하다/필요하다"의 정치가 아니라 "이번 릴리스가 몫의 몇 분을 먹는가"의 계산이 된다. 실습에서 수강생은 자기 시나리오의 가용성 목표를 분으로 환전하고 릴리스 몫을 나누어 본다 — 99.9\%면 43.8분, 릴리스 40\%면 17.5분, DORA 실측이면 월 0.6회다. 예상 소요 6분.}
\end{frame}

\begin{frame}{Incident와 Problem, Known Error DB — 복구로 닫는 것과 원인 제거로 닫는 것}
\begin{tbl}[\scriptsize]{L{0.1\textwidth}L{0.3\textwidth}L{0.36\textwidth}L{0.14\textwidth}}
KE & 증상 \ra{} 우회 & 근본 원인 · 상태 & 소유자 \\ \midrule
KE-01 & 구형 단말 9점 인쇄 8초 \ra{} 재부팅 & 펌웨어 호환 — 협의 중(이관 2번) & 오정근 \\
KE-02 & 폴백 배치 겹침(R-16) \ra{} 22:30 & 3PL 준실시간 전환 시(2027-06, 이관 4번) & 한동석·박준영 \\
KE-03 & 포인트 이중차감 \ra{} 중복 검출 배치 & 멱등 키 — 03-12 핫픽스 제거, 감시 잔류 & 수행사 \\
KE-04 & IF-SAP-12 타임아웃 \ra{} 수동 재전송 & ECC 야간 배치 경합 — EoM 2027-12까지 유지(이관 7번) & 박세연 \\
KE-05 & 구형 xls 수량 0 \ra{} xlsx 안내 & 사용률 1\% 미만 — 보류(CAB 04-14) & 박준영 \\
\end{tbl}
\keymsg{인시던트 140 · 문제 9 = 근본 제거 3 + Known Error 6 — 상태가 다르므로 소유자가 다르고, 소유자가 다르므로 이관처가 다르다}
\note{교재 3.7절, 워크북 §2 항목 4(KEDB). Incident Management는 복구, Problem Management는 원인 제거. 인시던트는 원인을 몰라도 복구되면 닫히고, 문제는 재발이 없어도 원인이 제거되어야 닫힌다. Known Error = 원인이 알려졌고 우회가 있으나 근본 제거가 아직인(또는 안 하기로 한) 오류. 표에 없는 KE-06(옴니 매장 간 재고 이동 예약 반영 15분 지연, KE-02와 동일 원인 계열, 오정근)까지 6건. 상태 네 종 — 제거 완료(KE-03), 다른 과제 의존(KE-02·06 \ra{} 3PL, KE-04 \ra{} ECC), 의도적 보류(KE-05), 협의 중(KE-01). Day3 등록부의 R-13(구형 단말)·R-16(배치 충돌)이 KE-01·KE-02가 된 것이 "잔여 리스크의 운영 이관"의 실물(4.6절, 3교시). "보류" 결정을 적지 않으면 6개월 뒤 "왜 안 고쳤나"가 된다. 예상 소요 6분.}
\end{frame}

\begin{frame}{개인정보 사건 — 조문은 단정, 적용은 비단정, 시한은 단정}
\fig[0.52]{../그림/PM_Day4/fig_3_6_privacy_timeline.pdf}
\figcap{그림 3-6. P1-02 포인트 이중차감(2,987명, 03-05 14:20 인지) — CPO 의견 20h40m \ra{} 결정 \ra{} 통지·복구 27h40m \ra{} 신고 49h40m \ra{} 시한 72h. 이 절차가 운영 문서에 없어 회의를 소집한 사실이 L-14}
\keymsg{개정법(제21445호, 2026-09-11 시행) — '유출등'에 훼손 포함·72시간은 단정, 이 사건이 훼손인지는 비단정, 그래도 시한 안에 결정한다}
\note{교재 3.8절, 워크북 §0.1 규칙 7·§2 항목 8. 조문 수준에서 단정 가능한 것: 유출등의 정의에 위조·변조·훼손 포함, 인지로부터 72시간 내 정보주체 통지·보호위원회 신고(대상 요건은 시행령·고시), 과징금 상향. 단정할 수 없는 것: ① 시스템 오류로 잔액이 틀리게 기록된 것이 '훼손'인가(시행 초기 해석례 부재), ② 2,987명이 규모 요건에 해당하는가, ③ 복구 후에도 신고 의무가 유지되는가. 최영주 CPO 서면 의견(03-06 11:00): 배제할 수 없으므로 보수적으로 72시간 내 신고 권고, 통지·복구는 규제 해당 여부와 무관하게 즉시. 결정 03-06 14:00(정미란·최영주·박준영·PM) — 통지·복구 03-06 18:00, 신고 03-07 16:00(49h40m), 협의회 서면 03-08. 기록을 "법상 신고 대상이므로 신고"로 정리하면 깔끔하지만 거짓이다 — 셋을 분리해 적어야 다음 사건에서 "지난번 쟁점 ①이 그 후 결정례로 정리되었는가"를 볼 수 있다. CPO의 의견은 검토 의견이지 판정이 아니다. 예상 소요 8분.}
\end{frame}

\begin{frame}{사용자 채택과 변화관리 — 저항은 정보다}
\begin{tbl}[\scriptsize]{L{0.18\textwidth}L{0.3\textwidth}L{0.12\textwidth}L{0.3\textwidth}}
집단 & 지표(2주 \ra{} 8주) & 장벽 & 조치 \\ \midrule
가맹 점주 388점 & 팩스·엑셀 발주 24\% \ra{} 11\% (목표 $\le$15) & K/A & P3 방문 교육 2회차 · R5 간편 재발주 \\
가맹 점주 & 마감 20:00 준수 81 \ra{} 96\% & R & 마감 30분 전 알림 \\
직영 227점 & 반품 절차 준수 78 \ra{} 94\% & \textbf{R} & POS 강제 확인 단계 + 지역매니저 점검 \\
서비스데스크 11명 & 1차 해결률 52 \ra{} 73\% & K & FAQ 20건 · KEDB 우회 공유 \\
물류센터 246명 & 재고 조회 사용률 87\%(4주) & A/D & 한동석 서면 04-12 \\
\end{tbl}
\keymsg{채택은 교육 이수율이 아니라 사용 지표로 잰다 — R 장벽에 교육을 더 하면 ELS가 끝나는 날 내려간다}
\note{교재 3.9절(그림 3-7), 워크북 §2 항목 9. ADKAR(Prosci, Hiatt 2006) — Awareness·Desire·Knowledge·Ability·Reinforcement, 어느 단계에 막혀 있는지의 진단이 먼저다. Kotter의 단기 성과 공지 — 04-05 협의회 공문 "발주 리드타임 38h \ra{} 13.6h". Bridges의 끝냄 — OD-POS 종료 세션 03-19. Ford·Ford·D'Amelio(2008): 저항은 극복할 장애가 아니라 변화의 설계에 관한 정보이고 때로는 추진자 자신이 만든 것이다 — 강윤지 "이번엔 없어지지 않는 건가요"는 2019년의 기억이 만든 D 장벽, 김미르 "처리량 목표가 인력 절반이라는 뜻이라면"은 편익 지표와 인력 평가의 연결에 대한 A 장벽. 가맹의 채택률이 오른 것은 교육이 아니라 부속서 동의(12-22, 91\%)로 D가 먼저 풀렸기 때문이고, 직영의 절차 미준수(R)가 P1-02의 노출면이었다 — 채택과 인시던트가 같은 자리에서 만난다. 예상 소요 6분.}
\end{frame}

\begin{frame}{D$-$2 라이선스 사건 — 거버넌스 변경의 부작용이 14개월 뒤 청구되는 구조}
\fig[0.55]{../그림/PM_Day4/fig_2_6_license_chain.pdf}
\figcap{그림 2-6. 내부통제 지침(2026-02-16) · 50/50 계약 변경(03-13) · 조달 상태 표 "계약 내" \ra{} 02-13 CFO 결재 대기 \ra{} 02-24 14:30 발견 \ra{} ② 임시 라이선스 0.02억 \ra{} L-09 \ra{} 지침 개정 요청 2·3}
\keymsg{세 개의 정상적인 결정이 겹쳐 만든 대기 — 담당자가 아니라 규정의 겹침을 적어야 L-09가 지침 개정이 된다}
\note{교재 2.8절(1교시에서 넘긴 절), 워크북 §1 항목 9(런북 부록). 2027-02-24(수) 14:30 컨트롤룸 사전 점검에서 API GW 운영 환경 2차분 발주서 "미접수" 회신. 발주서는 02-13 박세연 결재 후 CFO 직접 결재 대기 — 재무본부 내부통제 지침(2026-02-16)이 5,000만 원 이상 단일 건을 CFO 직접 결재로 바꿨고, 정미란은 출장 중. 선택지 ① 결재 대기 \ra{} 03-04, 창 이탈 / ② 벤더 임시 라이선스 30일 0.02억(컨틴전시 미배정 풀, CCB 03-18 사후 보고) / ③ 개발 키 운영 사용 — 계약 위반, 기각. 원인은 실수가 아니라 (1) R-04 대응으로 라이선스를 50/50으로 바꾼 계약 변경(IS-01, 2026-03-13)이 2차분 발주를 2027-02로 밀었고 (2) 그 시점이 지침 시행 후였으며 (3) 조달 상태 표가 "계약 내"로만 표기해 발주 상태가 보이지 않았다는 것. 그래서 런북 부록에 적었고 다음 컷오버 게이트 1 조건에 "라이선스·결재 상태"가 들어가며, L-09 \ra{} 내부통제 지침 개정 요청 2번(게이트·창 앞 4주 이내 발주 건 사전 일괄 결재). 이 사건이 2교시에 있는 이유 — 서비스 전환의 사건 대부분은 기술이 아니라 절차의 겹침에서 온다. 예상 소요 6분.}
\end{frame}

\begin{frame}{2교시 핵심 정리 — 약속과 그 약속을 받치는 것}
\begin{itemize}
\item ITIL 현행은 v5 · 이 교재는 practice 어휘로 고정 — "4가 최신"이라 쓰지 마라
\item SAC = 인수자 판정 + 헌장 거부권 + 조건부 경로 · 미해소 7 \ra{} 충족 11·조건부 1 \ra{} 이관 1번
\item ELS 8주 = 기간 · n주 연속 종료 기준 · 연동된 인력 축소(20\ra{}4)
\item SLA $\ge$ OLA·UC 사슬의 합 — 315 $>$ 120 재협상, 3PL UC 없음 강등
\item 99.5\% = 219.2분 · 릴리스 몫 87.7 · 복구 240분이면 월 3회, 45분이면 19.5회
\end{itemize}
\keymsg{인시던트는 복구로, 문제는 원인 제거로, 개인정보는 조문·적용·시한 세 줄로, 채택은 사용 지표로 — 워크북 §2 ⑰로 간다}
\note{교재 제3장 핵심 정리 5항. 4교시 실습 ①에서 워크북 §2 ⑰ 서비스 전환 패키지를 쓴다 — SAC 12항목의 판정자·증빙 열, ELS 종료 기준 다섯과 인력 축소 연동, 3층 매핑표에서 가장 긴 사슬 합산, error budget policy(환전 계산 포함), KEDB 6건의 상태·소유자·이관처. 수강생에게 미리 말할 것 둘 — 자기 시나리오의 SLA 항목을 하나 골라 3층을 적어 보면 대부분 UC가 없거나 UC가 SLA보다 길다는 것을 발견한다 / 강등은 실패가 아니라 재협상 안건이다. 3교시는 종료(제4장) — 종료 보고서·BL-02 대비 최종 성과·계약 종결·이관 10건. 예상 소요 4분.}
\end{frame}
